\chapter{TESTING}

\section{Testing Methodology}
Testing is a critical phase in the development of CivicGuard, ensuring that citizens can reliably report issues and that municipal officers can manage them effectively. We adopted a **V-Model testing strategy**, where each development phase is mirrored by a specific testing activity. This ensures that requirements are validated early and defects are minimized.

\section{Unit Testing}
Unit testing focuses on verifying the individual components and business logic of the application in isolation. We use **PHPUnit**, the world's most popular PHP testing framework, to execute these tests. By mocking external dependencies (like the filesystem or database), we can ensure that a specific method—such as the `ReportRoutingService`'s decision logic—works perfectly under every possible edge case. These tests are the first line of defense against logic bugs.

\renewcommand{\arraystretch}{1.3}
\begin{longtable}{|p{1.5cm}|p{2.5cm}|p{4cm}|p{4cm}|c|}
    \hline
    \textbf{Test ID} & \textbf{Component} & \textbf{Description} & \textbf{Expected Result} & \textbf{Status} \\
    \hline
    \endfirsthead
    \hline
    \textbf{Test ID} & \textbf{Component} & \textbf{Description} & \textbf{Expected Result} & \textbf{Status} \\
    \hline
    \endhead
    \hline
    \caption{Unit Test Cases for CivicGuard} \label{tab:unit_tests} \\
    \endlastfoot

    UT01 & Auth & Login with valid credentials & Redirect to dashboard & Pass \\
    \hline
    UT02 & Auth & Register new user & Account created in DB & Pass \\
    \hline
    UT03 & Form & Validate mandatory photo & Error message shown & Pass \\
    \hline
    UT04 & Dashboard & Filter user reports & Only own reports visible & Pass \\
    \hline
    UT05 & Search & Search by report ID & Specific report is returned & Pass \\
    \hline
    UT06 & Admin & Create Announcement & Broad-casted to all users & Pass \\
    \hline
    UT07 & Validation & Check character limit on description & Prevents submission if too long & Pass \\
    \hline
    UT08 & Logic & Test department assignment mapping & Correct ID assigned based on category & Pass \\
    \hline
\end{longtable}

\section{Integration Testing}
Integration testing ensures that the frontend Blade components, backend controllers, and the PostgreSQL storage layer interact as a cohesive unit. This phase is critical for verifying "Data Integrity"—ensuring that a report submitted through the browser actually ends up in the database with the correct binary image path and departmental foreign keys. It validates the "Handshake" between disparate system modules.

\begin{table}[H]
    \centering
    \renewcommand{\arraystretch}{1.3}
    \small
    \begin{tabular}{|c|p{3.5cm}|p{3.5cm}|p{3.5cm}|c|}
        \hline
        \textbf{Test ID} & \textbf{Component} & \textbf{Test Description} & \textbf{Expected Result} & \textbf{Status} \\
        \hline
        IT01 & Form to DB & Submit a complete report & Record appears in DB with correct values & Pass \\
        \hline
        IT02 & Storage & Upload an image asset & File exists in public storage directory & Pass \\
        \hline
        IT03 & Notification & Officer updates status & Notification sent to user's dashboard & Pass \\
        \hline
        IT04 & Admin Panel & Approve/Reject reports & Status change reflected in UI and database & Pass \\
        \hline
        IT05 & Relation & Delete Report & Associated notifications also deleted & Pass \\
        \hline
    \end{tabular}
    \caption{Integration Test Cases for CivicGuard}
\end{table}

\section{System \& Performance Testing}
System testing evaluates the end-to-end functionality of the entire CivicGuard ecosystem in a production-like setting. This includes verifying the session persistence across different tabs and the responsiveness of the dashboard on mobile devices. Performance testing, on the other hand, focuses on "Stress Testing"—observing how the server handles concurrent high-resolution photo uploads to ensure that the user experience remains smooth even during peak reporting hours.

\begin{table}[H]
    \centering
    \renewcommand{\arraystretch}{1.3}
    \begin{tabular}{|c|p{2.5cm}|p{4cm}|p{4cm}|c|}
        \hline
        \textbf{Test ID} & \textbf{Type} & \textbf{Test Description} & \textbf{Expected Result} & \textbf{Status} \\
        \hline
        ST01 & System & Full Life-cycle scenario & Report $\rightarrow$ Update $\rightarrow$ Resolve & Pass \\
        \hline
        ST02 & Access Control & Citizen tries to reach admin dash & Access denied & Pass \\
        \hline
        PT01 & Performance & Dashboard load time (100 reports) & Under 1.5 seconds & Pass \\
        \hline
        PT02 & Performance & Image stress test (multi-upload) & No server timeout & Pass \\
        \hline
    \end{tabular}
    \caption{System and Performance Testing Results}
\end{table}

